% ===========================================================================
% Chapter 1: Introduction
% ===========================================================================
\section{Introduction}

\subsection{Motivation}

Accurate crop classification from satellite imagery is fundamental to food security
monitoring, agricultural subsidy verification, and climate impact assessment.
Modern remote sensing platforms provide multi-modal observations: optical imagery
(Landsat 8/9, Sentinel-2), Synthetic Aperture Radar (Sentinel-1), and Digital
Elevation Models (SRTM, ALOS). Each modality captures complementary information:
optical reflectance reveals vegetation spectral signatures, SAR backscatter
penetrates cloud cover and senses canopy structure, and DEM provides topographic
context that modulates both.

However, fusing these modalities presents four fundamental mathematical
challenges that have not been adequately addressed by existing deep learning
approaches:

\begin{enumerate}
    \item \textbf{Terrain Geometry:} DEM features are typically encoded by standard
    CNNs as implicit $128$-dimensional vectors with no geometric semantics.
    Cross-region generalization fails because CNN features lack SE(3)-invariance,
    and autograd-based second-derivative computation causes computational graph
    explosion ($4.7$s, $12.4$GB VRAM for $1024^2$ at $5$ channels).

    \item \textbf{Cross-Domain Alignment:} Source domain (ImageNet/SeCo pre-training)
    and target domain (specific crop regions) exhibit distribution shift that
    simple channel-wise affine transformations cannot capture. Standard
    Gromov-Wasserstein distance computation is $O(N^3)$, prohibitive for
    per-batch domain adaptation.

    \item \textbf{Multi-Modal Evidence Conflict:} When optical, SAR, and DEM
    modalities disagree (e.g., optical says ``wheat'' while SAR says ``rice''),
    standard Dempster-Shafer combination can produce the Zadeh paradox
    (high-confidence wrong answer). Scalar vacuity cannot distinguish random
    noise from structural semantic contradiction.

    \item \textbf{Temporal Generalization:} Dynamic convolution layers are
    empirically effective but lack rigorous mathematical justification,
    making architecture design and hyperparameter selection dependent on
    expensive grid search ($\sim12$ hours).
\end{enumerate}

\subsection{Contributions}

This work makes the following contributions:

\begin{enumerate}
    \item \textbf{Differential Geometric DEM Encoding (Module 1):} We replace
    CNN-based DEM encoding with SIREN implicit surface fitting and closed-form
    computation of five geometric invariants via the method of moving frames,
    proving SE(3)-invariance (Theorem 1) and isometry-invariance (Theorem 2).

    \item \textbf{Joint Optimal Transport--Grassmann Alignment (Module 2):} We
    formulate cross-modal feature alignment as a unified optimization over
    sliced Gromov-Wasserstein distance and Grassmann geodesic distance, solved
    via Stiefel ADMM with closed-form V-subproblem (Theorem 4) and linear
    convergence guarantee.

    \item \textbf{Topological Evidence Fusion (Module 3):} We prove the DS-Cup
    Product Equivalence Theorem, establishing that Dempster-Shafer combination
    is strictly equivalent to the cup product on simplicial cochain complexes.
    We show that conflict between modalities is a cohomological obstruction
    (Theorem 2) and introduce three-way conflict classification with persistent
    homology correction.

    \item \textbf{Statistical Learning Guarantees (Module 4):} We prove that
    dynamic convolution operators are equivalent to kernel-space low-rank linear
    attention (Theorem 1) and derive temporal Rademacher generalization bounds
    under geometric mixing conditions (Theorem 2), yielding closed-form optimal
    hyperparameters.

    \item \textbf{FusionCropNet V6 Architecture:} An integrated model with 14
    components implementing all four mathematical modules, achieving $313$
    passing tests, $48.9$M parameters, and CPU-trainability.
\end{enumerate}

\subsection{Related Work}

\subsubsection{Deep Learning for Crop Classification}

Deep learning has become the dominant paradigm for remote sensing crop
classification. Early works applied 2D CNNs (ResNet, DenseNet) to single-temporal
imagery~\cite{zhong2019deep, russwurm2018multi}. The introduction of temporal
attention mechanisms~\cite{russwurm2020breizhcrops, garnot2020satellite} enabled
multi-temporal fusion, but these methods treat each pixel independently,
discarding spatial context. 3D CNNs~\cite{ji2013convolutional} and
ConvLSTM~\cite{shi2015convolutional} partially address this but at prohibitive
computational cost for large-area mapping.

The Transformer architecture~\cite{vaswani2017attention} was adapted to remote
sensing via positional encodings of spatial-temporal coordinates, achieving
state-of-the-art performance on crop classification benchmarks. However,
Transformer-based methods suffer from $O(T^2)$ complexity in the temporal
dimension and lack principled treatment of multi-modal uncertainty.

\subsubsection{Mathematical Foundations in Deep Learning}

\textbf{Differential geometry} has been applied to deep learning primarily
through manifold learning and geometric deep learning~\cite{bronstein2017geometric}.
SIREN~\cite{sitzmann2020siren} introduced periodic activation functions for
implicit neural representations, but its application to DEM encoding and
closed-form geometric invariant extraction is novel to this work.

\textbf{Optimal transport} has gained traction in domain adaptation through
Wasserstein distance minimization~\cite{courty2017optimal}. The Gromov-Wasserstein
distance~\cite{peyre2016gromov} extends OT to metric measure spaces, and sliced
variants~\cite{bonneel2015sliced} reduce computational complexity. Joint
OT-subspace alignment on the Grassmann manifold~\cite{edelman1998geometry} is
a contribution unique to this work.

\textbf{Algebraic topology} has seen limited application in machine learning,
primarily through persistent homology for data analysis~\cite{carlsson2008topology}
and topological data analysis (TDA)~\cite{edelsbrunner2008persistent}. The
connection between Dempster-Shafer theory and cup products on simplicial
cochain complexes is, to our knowledge, first established in this work.

\textbf{Statistical learning theory} provides generalization guarantees
through Rademacher complexity~\cite{bartlett2002rademacher} and PAC-Bayesian
bounds~\cite{mcallester1999pac}. The temporal generalization bound under
geometric $\beta$-mixing extends these results to non-i.i.d. time series
data~\cite{mohri2018foundations, kuznetsov2017generalization}.

\subsubsection{Evidential Deep Learning}

EDL~\cite{sensoy2018evidential} models prediction uncertainty through
Dirichlet distributions, providing vacuity and dissonance estimates.
Subsequent works extended EDL to classification~\cite{malinin2018predictive}
and regression tasks. However, existing EDL methods treat vacuity as a scalar,
failing to capture the \textit{structure} of multi-modal conflict. Our
topological reformulation addresses this fundamental limitation.

\subsubsection{Multi-Modal Remote Sensing Fusion}

Multi-modal fusion in remote sensing typically employs early fusion
(concatenation), late fusion (ensemble averaging), or attention-based
intermediate fusion~\cite{schmitt2020sen12ms}. Self-supervised pre-training
methods such as SeCo~\cite{manen2022seco} provide remote sensing-specific
representations but do not address cross-modal alignment geometry.
FusionCropNet V6 advances beyond these methods by (1) establishing
mathematical foundations for each fusion mechanism, (2) providing
SE(3)-invariant geometric features, and (3) detecting structural
cross-modal conflict through algebraic topology.

\subsection{Paper Organization}

Section 2 presents the four mathematical theory modules with complete formal
theorem statements, proofs, and cross-module synergy derivations. Section 3
describes the V6 architecture with forward-propagation pseudocode and the
unified loss function. Section 3B provides geographic and remote sensing
context. Section 4 reports experimental validation with training protocols.
Section 5 concludes with limitations and future work. The Appendix contains
code listings and Lean 4 formalization details.
